% --- Данные титульного блока; вёрстку выполняет \maketitle ---
% (переопределение \@maketitle — в inc/preamble.tex)

\udk{XXX.XX}

\title{СРАВНИТЕЛЬНЫЙ АНАЛИЗ ЛИНЕЙНОЙ (LSA) И НЕЛИНЕЙНОЙ
(ALL-MINILM-L6-V2) МОДЕЛЕЙ ВЕКТОРНОГО ПРЕДСТАВЛЕНИЯ ТЕКСТА}
\author{\textsuperscript{1}Шапаренко Ф.\,А.}
\affiliation{\textsuperscript{1}МИРЭА~--- Российский технологический
университет, Москва}
\email{shaparenko.fedor@gmail.com}

\abstractru{В работе проводится сравнительный анализ двух подходов к
  векторному представлению текста: линейной модели на основе
  латентно-семантического анализа (LSA) и нелинейной модели-трансформера
  \texttt{all-MiniLM-L6-v2}. Цель~--- выявить связь между алгебраическими
  свойствами векторных пространств моделей и качеством их работы на
  прикладных задачах. Сравнение проведено на датасетах STS Benchmark, 20
  Newsgroups и MRPC с использованием интринсивных (participation ratio,
  effective rank, isotropy score, hubness score) и экстринсивных (корреляция
  Спирмена, Accuracy, F1-macro, ARI, silhouette) метрик. Показано, что
  нелинейная модель существенно превосходит LSA в задачах, требующих учёта
  контекста (семантическое сходство и кластеризация), тогда как на задачах,
  опирающихся на совпадение ключевых слов (тематическая классификация и
  обнаружение парафраз), модели сопоставимы. Установлено, что высокие
  значения геометрических характеристик пространства (PR, ER, изотропия) не
являются достаточным условием семантического качества представлений.}
\keywordsru{латентно-семантический анализ, векторное представление текста,
  эмбеддинги, трансформеры, all-MiniLM-L6-v2, сингулярное разложение,
семантическое сходство, изотропия пространства, кластеризация текстов.}

\titleen{COMPARATIVE ANALYSIS OF LINEAR (LSA) AND NONLINEAR
(ALL-MINILM-L6-V2) TEXT VECTOR REPRESENTATION MODELS}
\authoren{\textsuperscript{1}Shaparenko F.\,A.}
\affiliationen{\textsuperscript{1}MIREA~--- Russian Technological University,
Moscow}
\abstracten{This paper presents a comparative analysis of two approaches to
  text vector representation: a linear model based on latent semantic analysis
  (LSA) and the nonlinear transformer model \texttt{all-MiniLM-L6-v2}. The aim
  is to reveal the relationship between the algebraic properties of the
  models' vector spaces and their performance on applied tasks. The comparison
  is carried out on the STS Benchmark, 20 Newsgroups, and MRPC datasets using
  intrinsic (participation ratio, effective rank, isotropy score, hubness
  score) and extrinsic (Spearman correlation, Accuracy, F1-macro, ARI,
  silhouette) metrics. The results show that the nonlinear model substantially
  outperforms LSA on context-dependent tasks (semantic similarity and
  clustering), whereas on tasks driven by keyword overlap (topic
  classification and paraphrase detection) the models are comparable. High
  values of geometric space characteristics (PR, ER, isotropy) are shown to be
insufficient as a condition for the semantic quality of representations.}
\keywordsen{latent semantic analysis, text vector representation, embeddings,
  transformers, all-MiniLM-L6-v2, singular value decomposition, semantic
similarity, space isotropy, text clustering.}

\maketitle
